\section{Introduction}
\label{sec:introduction}

Large language models have become increasingly capable of generating executable source code from natural-language specifications, making code generation one of the most prominent applications of generative AI in software engineering. However, plausible-looking code is not necessarily correct. Generated programs may violate the syntax of the programming language, fail during execution, or execute normally while deviating from the intended functionality. These failures are increasingly discussed as \emph{code hallucinations}: outputs that appear plausible but are incorrect or inconsistent with the requested behavior~\citep{pavel2025hallucination,lee2025hallucination}. \citet{lee2025hallucination}, for example, distinguish syntactic, runtime-execution, functional-correctness, and code-quality hallucinations. In this study, we focus on the first three categories, which correspond to observable failures in correctness and executability. Understanding and detecting these failures is important because generated code can appear convincing enough to be accepted by a developer even when its behavior is incorrect, transferring part of the verification burden from code generation to code assessment.

Detecting such hallucinations is itself a difficult problem. Syntax errors can often be identified mechanically, but executable code may fail only for particular inputs or silently return incorrect results. Empirical studies have shown that functional failures represent an important fraction of errors produced by code-generation models, while different models exhibit substantially different failure distributions~\citep{dou2025whatswrong,tian2025codehalu}. Test execution can provide an objective oracle when sufficiently strong test cases are available, but a model evaluating generated code does not necessarily have access to those tests or to execution results. Moreover, generating correct code and recognizing incorrect code are distinct capabilities: a model that performs well as a generator may still be a poor evaluator of its own or another model's output. This motivates studying hallucination detection as a task in its own right rather than treating correctness assessment only as a by-product of code generation.

We investigate this problem from the perspective of \emph{Small Language Models} (SLMs). Our central idea is to evaluate whether compact models can not only generate code, but also act as hallucination detectors, and whether task specialization can compensate for limited model scale. Rather than reducing verification to a binary correct/incorrect decision, we formulate hallucination detection as a multilabel classification task over \emph{Syntax}, \emph{Runtime}, \emph{Functional}, and \emph{Correct}. A generated solution may contain both runtime and functional failures when different test cases expose different behaviors, while \emph{Correct} is reserved for solutions that successfully satisfy the complete evaluation suite. This formulation allows us to study not only whether an SLM recognizes that a program is wrong, but also which forms of hallucination it is capable of identifying. We further examine whether models can recognize errors in code that they generated themselves and whether a specialized SLM can provide a more reliable detector than general-purpose models.

Our study follows a controlled experimental pipeline built on the 563 programming tasks of EvalPlus~\citep{liu2023evalplus}. Multiple SLMs first generate one solution for each task, and the resulting programs are compiled and executed against the strengthened HumanEval and MBPP test suites to construct an execution-based multilabel ground truth. The same task descriptions and generated programs are then presented to the models without test results or reference labels, turning hallucination detection into an independent evaluation task. To investigate specialization, we construct a supervised dataset through knowledge distillation: three teacher models independently receive the execution-derived labels and produce natural-language rationales explaining the corresponding failures. These rationales are used as auxiliary supervision when fine-tuning Gemma-based models with QLoRA, while the quantitative evaluation remains based solely on the predicted classes. We evaluate different optimization strategies to account for the strong class imbalance in the dataset and compare the resulting specialized detectors with both general-purpose SLMs and substantially larger models under the same detection protocol. Finally, we investigate whether conventional static-analysis signals from Python compilation and Pyright can complement model-based generation and hallucination detection.

The results reveal a clear separation between code-generation capability and hallucination-detection capability. Across the evaluated SLMs, code-generation correctness ranges from 62\% to 81\%, with functional hallucinations being considerably more frequent than syntax failures. Error detection is substantially more heterogeneous: the best general-purpose SLM reaches 73.1\% exact-set Accuracy and 53.8\% Macro-F1, while no general-purpose model detects more than approximately one third of the Runtime errors. Models also differ substantially in their ability to recognize their own failures; self-error detection ranges from only 5.3\% for GPT-OSS 20B to 69.3\% for Qwen 3.8 27B, showing that strong generation does not imply strong self-evaluation. Specialized training substantially improves detection: the Accuracy-oriented model reaches 77.8\% Accuracy, a gain of 31.9 percentage points over its base model, while Macro-F1-oriented training reaches 54.2\% Macro-F1 and provides a more balanced error profile. Moreover, the specialized SLMs remain competitive with models that are an order of magnitude larger: the best specialized detector surpasses both evaluated large models in Macro-F1, while the Accuracy-oriented detector also achieves the highest Accuracy. These results indicate that task specialization can partially compensate for model scale, although Runtime hallucinations remain challenging across all evaluated model groups.

% TODO RQ5: After the tool-assisted experiments are finalized, add one or two sentences here summarizing the strongest result regarding (i) static analysis alone, (ii) tool-assisted code generation, and/or (iii) tool-assisted hallucination detection. Prefer a concrete improvement/regression value rather than a qualitative statement.

The main contributions of this work are:
\begin{itemize}
    \item An empirical evaluation of the code-generation and hallucination-detection capabilities of SLMs under a common execution-based ground truth, distinguishing syntactic, runtime, and functional failures through multilabel classification.
    \item An analysis of \emph{self-error detection}, measuring whether models can recognize the hallucinations present in code that they generated themselves and showing that generation and self-evaluation capabilities do not necessarily coincide.
    \item A specialized SLM-based hallucination detector trained with execution-derived labels and teacher-generated rationales, together with an analysis of how different optimization objectives and class-imbalance strategies affect its detection behavior.
    \item A controlled comparison between general-purpose SLMs, specialized SLMs, and substantially larger language models under the same detection protocol, showing that increased scale does not guarantee superior hallucination detection.
    \item An investigation of the complementarity between model-based reasoning and traditional static-analysis evidence, considering static tools both as feedback during code generation and as supporting evidence during hallucination detection.
\end{itemize}

The remainder of this paper is organized as follows. Section~\ref{sec:related-work} reviews existing work on code hallucinations, empirical failure characterization, automatic code evaluation, specialized critics, and feedback-guided code generation. Section~\ref{sec:methodology} describes the evaluated models, generation protocol, execution-based ground truth, hallucination-detection task, knowledge-distillation procedure, specialized training, evaluation metrics, and integration with static-analysis tools. The subsequent sections present the empirical results for code generation, general hallucination detection and self-evaluation, specialized detection, comparison with larger models, and tool-assisted approaches. Finally, the paper discusses the implications and limitations of the findings and summarizes the main conclusions.
